\section{Scope: complementing the middle atlas without asserting existence}
The preceding middle-channel result proves a complete two-chart atlas with a
square-root cutoff \cite{Middle}. It does not prove that that atlas is occupied.
The present continuation treats the other channel. It gives a three-chart
exterior atlas indexed by an explicit cubic cutoff, preserves the original
square-divisor availability in the norm chart, and disproves transferring the
middle square-root bound unchanged to the exterior cofactors. The universal
existence step remains unproved. In particular, a finite decision algorithm is
not a witness-guaranteed algorithm.

The exterior normalization below is an explicitly marked presentation of the
classical Type~I parametrization; compare \cite[\S2, Proposition~2.2]{ET} and
the recent divisor-parametrization comparison \cite[\S3]{BHF}.  The paired
middle channel is Type~II in the sense of \cite[\S2,
Proposition~2.6]{ET}.  Every algebraic identity used here is proved below in
the retained marking.  No historical-priority claim is made for the derived
cutoff or atlas.

\section{The original exterior state and its three arithmetic parameters}
Let $p,a,u\in\mathbb Z_{>0}$, with $p>1$ and $p\equiv1\pmod4$, and retain
the unit first-half source
\[
 p/4<a<p/2,\qquad \gcd(p,a)=1,\qquad R=4a-p,
 \qquad u\in\mathbb Z_{>0},\quad u\mid a^2,\quad R\mid4u+1. \tag{1}
\]
The cutoff theorem applies even when this integer parameter $p$ is composite.
When asserting a complete prime ES atlas, primality is imposed separately.
Write $a=\prod_\ell\ell^{e_\ell}$ and retain every original exponent
$\beta_\ell=v_\ell(u)-e_\ell\in[-e_\ell,e_\ell]$.
Set
\[
 g=\gcd(a,u),\quad
 (h,r,s)=(g^2/u,u/g,a/g),\quad
 \kappa=(pr+s)/R.                                        \tag{2}
\]
Prime valuations prove that these are positive integers, $a=hrs$, $u=hr^2$,
and $\gcd(r,s)=1$. The gate in (1), after multiplication by the unit $s$ modulo
$R$, gives $R\mid pr+s$. The ordered return is
\[
 \boxed{(x,y,z)=(a,hs\kappa,phr\kappa),\qquad
 \frac4p=\frac1x+\frac1y+\frac1z.}                       \tag{3}
\]
It is checked by multiplication. Its ordered inverse is
\[
 u=\frac{a^2}{Ry-pa}.
\]
No denominator permutation is suppressed.

Retain the complement divisor and the exterior cofactor
\[
 \boxed{v=\frac{a^2}{u}=hs^2,\qquad
 D=\frac{4u+1}{R}=\frac{r+\kappa}{s}.}                   \tag{4}
\]
The second equality follows from
$RsD=s(4hr^2+1)=(p+R)r+s=R(r+\kappa)$.
Consequently
\[
 \boxed{pD+1=4hr\kappa,\qquad
 (p+R)^2+4v=4vRD.}                                     \tag{5}
\]
Here $R,D\equiv3\pmod4$, and $\kappa>r$ because $R<p$ and
$R\kappa=pr+s$. Thus
\[
 \boxed{p+R<2vD.}                                       \tag{6}
\]
The case $h=1$ would make $-1$ a square modulo $R\equiv3\pmod4$ through
$R\mid4r^2+1$.  Indeed, some prime $q\equiv3\pmod4$ occurs to odd
valuation in $R$; the divisibility also gives $q\nmid2r$ and
$(2r)^2\equiv-1\pmod q$, contradicting Euler's criterion.  Hence $h>1$.
Further,
$\gcd(h,R)=\gcd(h,D)=1$ by (1) and (5), so
\[
 \boxed{v\ne R,\qquad v\ne D.}                         \tag{7}
\]
This does \emph{not} assert $\gcd(v,D)=1$: the square factor $s^2$ may overlap
$D$. For example, at $p=37$, $(a,u,h,r,s,\kappa)=(12,8,2,2,3,7)$ has
$(v,D)=(18,3)$.

\section{A strict cubic cutoff}
\begin{theorem}[Exterior three-parameter bound]\label{thm:cubic}
For every state (1), put $w=\min(v,R,D)$. Then
\[
 \boxed{(p+w)^2>4(w+1)^2(w-1).}                         \tag{8}
\]
Let $C(p)$ be the largest positive integer $w$ satisfying (8). Then
\[
 1\le C(p)<p,\qquad \min(v,R,D)\le C(p),
 \qquad C(p)=(p/2)^{2/3}+O(p^{1/3}).                    \tag{9}
\]
The constant and strict sign in (8) are part of the statement. The asymptotic
notation is not used in the verifier: $C(p)$ is computed by exact integer
comparison.
\end{theorem}
\begin{proof}
Put $t=\min(R,D)$. If $R=t$, (5) gives
\[
 (p+t)^2\ge4v(t^2-1).
\]
If $D=t\le R$, direct subtraction using (5) gives
\[
 (p+t)^2-4v(t^2-1)
 =(R-t)(4vt-2p-R-t)\ge0.                                \tag{10}
\]
Indeed (6) implies $4vt-2p-2R>0$; the last factor in (10) is that positive
quantity plus $R-t$. Equality is possible only when $R=D=t$. Hence
\[
 p\ge 2\sqrt{v(t^2-1)}-t.                               \tag{11}
\]
By (7), either $v\ge t+1$ or $v\le t-1$.

In the first case $w=t$, and (11) gives
$p\ge2(w+1)\sqrt{w-1}-w$.
Equality would require $R=D=w$, $v=w+1$, and
$p+w=2(w+1)\sqrt{w-1}$. Thus $w-1$ would be an integer square, whereas
$w=R\equiv3\pmod4$ makes $w-1\equiv2\pmod4$. Equality is impossible.

In the second case $w=v$ and $t\ge w+1$. The function
$2\sqrt{w(t^2-1)}-t$ is strictly increasing for $t>1$, so
\[
 p\ge2w\sqrt{w+2}-w-1.
\]
This is strictly larger than $2(w+1)\sqrt{w-1}-w$. To verify the comparison,
put $U=w\sqrt{w+2}$ and $V=(w+1)\sqrt{w-1}$. For $w\ge2$,
\[
 U^2-V^2=w^2+w+1,\qquad U+V\le2w(w+1),
\]
so $U-V>1/2$. This proves the strict inequality in both cases; squaring its
positive sides proves (8).

The function
\[
 F(w)=2(w+1)\sqrt{w-1}-w
\]
is strictly increasing for $w>1$.
At $w=1$, (8) holds, and at $w=p$ it fails since
$(p+1)^2(p-1)-p^2=p^3-p-1>0$. There is therefore one final true integer in
$[1,p-1]$, obtained by binary search.  The defining inequality is precisely
$F(w)<p$.  Moreover
\[
 F(w)=2w^{3/2}-w+O(\sqrt w)\qquad(w\longrightarrow\infty).
\]
Maximality gives $F(C(p))<p\le F(C(p)+1)$, and the displayed expansion at
$C(p)$ and $C(p)+1$ gives
\[
 p=2C(p)^{3/2}+O(C(p)).
\]
In particular $C(p)\asymp p^{2/3}$; dividing by $2C(p)^{3/2}$ and taking the
$2/3$ power then yields
\[
 (p/2)^{2/3}=C(p)+O(C(p)^{1/2})
             =C(p)+O(p^{1/3}),
\]
which is (9) with the stated leading constant.
\end{proof}

The residue hypothesis in the strict form is essential.  If it is removed,
then
\[
 (p,a,u,R,D,v,w)=(19,6,6,5,5,6,5)
\]
satisfies the positive first-half divisor data but gives equality rather than
strict inequality:
\[
 (p+w)^2=24^2=576=4\cdot6^2\cdot4
              =4(w+1)^2(w-1).
\]
This is a boundary counterexample to the hypothesis-free statement, not an
Erd\H{o}s--Straus counterexample.

For a prime $p$, the least-nonresidue lower bound from the preceding work gives
$v=hs^2\ge h\ge\nu_p$. This may prune the complement chart, but it does not
replace the requirement that that chart actually contain a divisor satisfying
its congruence and availability condition.

\section{Three disjoint charts with complete inverses}
Define the root-divisibility kernel
\[
 K(v)=\prod_{\ell^e\parallel v}\ell^{\lceil e/2\rceil}.
\]
Then $v\mid a^2$ is equivalent to $K(v)\mid a$, prime by prime.
For the following atlas let $p\equiv1\pmod4$ be prime and $C=C(p)$.

\begin{theorem}[Complete original exterior atlas]\label{thm:atlas}
Every original first-half exterior state belongs to exactly one of the following
charts, and every chart entry reconstructs such a state.

\textbf{Direct residual chart.} Retain
\[
 3\le R\le C,\quad R\equiv3\pmod4,\quad a=(p+R)/4,
 \quad u\in\mathbb Z_{>0},\quad u\mid a^2,\quad R\mid4u+1.
\]
Compute $D=(4u+1)/R$, $v=a^2/u$, and retain $R\le D$, $R<v$.
The map and inverse are (1)--(4).

\textbf{Reciprocal chart.} Retain
\[
 \begin{gathered}
 3\le D\le C,\quad D\equiv3\pmod4,\quad B_D=(pD+1)/4,\\
 w\in\mathbb Z_{>0},\quad w\mid B_D^2,\quad
 w<B_D,\quad D\mid w+B_D.
 \end{gathered}
\]
Set
\[
 g=\gcd(B_D,w),\quad
 h=g^2/w,\quad r=w/g,\quad\kappa=B_D/g,
 \quad s=(r+\kappa)/D,\quad a=hrs,\quad u=w.
\]
Compute $R=4a-p$, $v=hs^2$, and retain $D<R$, $D<v$.

\textbf{Complement--norm chart.} Retain
\[
 \boxed{2\le v\le C,\quad R\mid p^2+4v,\quad
 R\equiv3\pmod4,\quad v<R<p,\quad
 R\equiv-p\pmod{4K(v)}.}                               \tag{12}
\]
Put
\[
 \boxed{a=(p+R)/4,\qquad u=a^2/v,\qquad
 D=(4u+1)/R,}                                           \tag{13}
\]
and retain $v<D$. Formula (2) then gives the complete original marking and (3)
its ordered denominators. The inverse reads the original complement $v=a^2/u$
and residual $R$.
\end{theorem}
\begin{proof}
For the direct chart, $p+R\equiv0\pmod4$ makes $a=(p+R)/4$ integral;
$0<R\le C<p$ gives $p/4<a<p/2$; and primality of $p$ gives
$\gcd(p,a)=1$.  The two displayed divisor conditions are exactly the
remaining source conditions, so
\[
 (R,u)\longmapsto ((p+R)/4,u),\qquad
 (a,u)\longmapsto(4a-p,u)
\]
are mutually inverse on the retained direct region.  In particular the full
$u$-fibre is retained; the outer residual $R$ alone is not treated as a state.

For the reciprocal chart put $B=B_D$ and
$g=\gcd(B,w)$.  Prime valuations in $w\mid B^2$ prove that
\[
 h=\frac{g^2}{w}\in\mathbb Z_{>0},\qquad
 w=hr^2,\qquad B=hr\kappa,\qquad \gcd(r,\kappa)=1.
\]
The inequality $w<B$ is exactly $r<\kappa$.  Since $pD+1=4B$, the odd
integer $D$ is coprime to $B$, hence to $g=hr$.  The condition
$D\mid w+B=g(r+\kappa)$ therefore gives $D\mid r+\kappa$, so
$s=(r+\kappa)/D$ is a positive integer.  If a prime divided both $r$ and
$s$, the identity $Ds=r+\kappa$ would make it divide $\kappa$, contrary to
$\gcd(r,\kappa)=1$; thus $\gcd(r,s)=1$.

Now
\[
 R\kappa=(4hrs-p)\kappa
 =s(4hr\kappa)-p\kappa
 =s(pD+1)-p\kappa=pr+s>0,
\]
and
\[
 DRs=R(r+\kappa)=Rr+R\kappa
 =(R+p)r+s=s(4hr^2+1).
\]
Cancelling the positive integer $s$ proves $DR=4u+1$ for
$u=w=hr^2$.  Furthermore
\[
 D(p-2a)=pD-2hr(r+\kappa)
 =4hr\kappa-1-2hr(r+\kappa)
 =2hr(\kappa-r)-1>0.
\]
Thus $R>0$, $p/4<a<p/2$, and, because $p$ is prime and $0<a<p$,
$\gcd(p,a)=1$.  We have recovered every condition in (1), not merely the
ordered reciprocal identity.  Conversely, for an original state in this
chart, $B_D=hr\kappa$ and $w=u=hr^2$, while
\[
 d\mid\gcd(r,\kappa)\Longrightarrow d\mid(R\kappa-pr)=s
 \Longrightarrow d=1,
\]
because $\gcd(r,s)=1$.  Therefore
\[
 \gcd(B_D,w)=hr\gcd(\kappa,r)=hr=g.
\]
The displayed formulas therefore recover $h,r,\kappa,s,a,u,R,D,v$ exactly.
The complementary reciprocal divisor would give $\kappa<r$ and is not silently
identified with this orientation.  All divisors $w$ in this chart are positive;
admitting signed divisors would destroy the reconstructed positive state.

For (12), let $d\mid\gcd(R,4v)$.  The norm divisibility
$R\mid p^2+4v$ gives $d\mid p^2$.  Since $p$ is prime and $0<R<p$, the
prime $p$ cannot divide $d$, and hence $d=1$.  Thus
$\gcd(R,4v)=1$. The residue modulo $4K(v)$ gives $K(v)\mid a$, equivalently
the actual integer availability $v\mid a^2$, so (13) defines
$u\in\mathbb Z_{>0}$. Finally
\[
 4v(4u+1)=(p+R)^2+4v\equiv0\pmod R.
\]
Cancelling the proved unit $4v$ gives the exterior gate and makes
$D=(4u+1)/R$ a positive integer. The range gives $p/4<a<p/2$; since $p$ is
prime it also gives $\gcd(p,a)=1$. Thus every condition in (1) has been
reconstructed, and gcd normalization applies. Conversely every original
state has $R\mid p^2+4v$ by (5), and its $v\mid a^2$ gives exactly the residue
modulo $4K(v)$.

The smallest of $R,D,v$ is at most $C$ by Theorem~\ref{thm:cubic}.
The grade $v$ cannot tie either of the other parameters.  If $R\le D$, then
either $R<v$, which is the direct chart, or $v<R\le D$, which is the norm
chart.  If $D<R$, then either $D<v$, which is the reciprocal chart, or
$v<D<R$, which is the norm chart.  Thus the three rules are disjoint and
exhaustive.  More precisely, an $R=D$ tie for the minimum, $R=D<v$, is
assigned to the direct chart, while $v<R=D$ belongs to the complement--norm
chart.  Both cases occur: $(p,a,u,R,D,v)=(13,4,2,3,3,8)$ is direct, whereas
$(5,2,2,3,3,2)$ is norm.
\end{proof}

Primality is not cosmetic in the generative atlas.  If it is dropped without
adding replacement coprimality conditions, the direct data
\[
 (p,R,a,u,D,v)=(21,3,6,2,3,18)
\]
pass the displayed arithmetic tests but have $\gcd(p,a)=3$; the reciprocal
data $(p,D,B_D,w)=(93,3,70,20)$ reconstruct $a=30$ with the same defect; and
the norm data $(p,v,R)=(25,5,15)$ fail the cancellation needed for
$R\mid4u+1$.  The cutoff and the classification of already-valid source
states do not use primality, but these three parameter-side converse maps do.

\paragraph{A correction to an overly strong coprimality shortcut.}
One may cancel $R$ against $v$ in (12), because it follows there from the
prime parameter and the norm divisibility. One may not assert
$\gcd(v,D)=1$ in the original state. The $p=37$ example above has a common
factor three. The reciprocal construction retains that factor through $s$.

\paragraph{Candidate size and the surviving existence question.}
There are at most $(C+1)/4$ values in each of the first two lists and $C-1$
values of $v$ in the third. Exact finite candidate counts are obtained by summing
$\tau(a^2)$, $(\tau(B_D^2)-1)/2$, and $\tau(p^2+4v)$ over their respective
lists, then applying the displayed original conditions. These are factorization
and divisor-enumeration bounds, not polynomial-time claims. The preceding
middle two-chart atlas has square-root-sized outer lists. Together the two
results decide the complete prime ES source at each given $p$. They still do
not prove that its result is nonempty for every $p$.

\section{An actual hard prime and an asymptotically extremal integer family}
At
\[
 \boxed{p=48049\equiv169\pmod{840}},
\]
complete trial division through $219$ proves primality. The original exterior
state is
\[
 \boxed{(a,u,R,D,v,h,r,s,\kappa)
 =(12090,24180,311,311,6045,6045,2,1,309).}                \tag{14}
\]
It gives
\[
 \boxed{\frac4{48049}=\frac1{12090}+\frac1{1867905}
 +\frac1{179501934690}.}
\]
Here $4u+1=311^2$, so the exterior cofactors genuinely share all their factors.
The preceding hard-prime middle cutoff is $247$. Both $R$ and $D$ exceed it.
Thus the middle square-root bound cannot be transferred unchanged to these
exterior parameters; the surviving map is Theorem~\ref{thm:atlas}, not a false
identification of the two channels.

For every integer $n\ge2$, set
\[
 h=4n^2-2,\quad r=n,\quad s=1,\quad
 R=D=4n^2-1,\quad\kappa=4n^2-n-1,
\]
\[
 p=16n^3-4n^2-8n+1,\qquad a=hn,\qquad u=hn^2.
                                                               \tag{15}
\]
Then $R^2=4hn^2+1$, $p=4a-R$, and $R\kappa=pn+1$.
Also $\gcd(p,a)=1$ since $\gcd(R,h)=\gcd(R,n)=1$.
The strict first-half inequalities and the ordered identity follow directly.
Thus these are actual unit integer sources, not arbitrary real triples satisfying
an inequality. Here
\[
 \min(v,R,D)=h=4n^2-2,
 \qquad \frac{h}{(p/2)^{2/3}}\longrightarrow1.
\]
In fact $C(p)=h$ at every member of (15). Writing $w=h$, its parameter is
$p=2w\sqrt{w+2}-w-1$, while
\[
 2(w+2)\sqrt w-w-1>2w\sqrt{w+2}-w-1=p.
\]
The left side is the bounding function at $w+1$. Thus (8) holds at $h$ and fails
at $h+1$. The integer cutoff itself is attained, including the prime examples
$p=97,929$.
The exponent $2/3$ and leading coefficient in the general integer-domain cutoff
cannot be decreased uniformly for these three parameters. At $n=2,4$, the
parameters $p=97,929$ are prime.  The complete congruence classification is
\[
 p(n)\equiv289\pmod{840}
 \quad\Longleftrightarrow\quad
 n\equiv24\ \hbox{or}\ 164\pmod{210}.                 \tag{16}
\]
Indeed reduction modulo $8,3,5,7$ gives respectively
$n\equiv0\pmod2$, $n\equiv0,2\pmod3$, $n\equiv4\pmod5$, and
$n\equiv3\pmod7$; the Chinese remainder theorem gives exactly the two
classes in (16).  No infinitude of prime values of the cubic is asserted;
consequently the asymptotic sharpness claim is for the stated unit integer
domain, not for a proved infinite prime subfamily.

The least complement grade is not automatically occupied. At $p=2521$, $v=11$
requires $R\mid6355485$ and $R\equiv31\pmod{44}$, $0<R<p$.
The factorization $6355485=3^2\cdot5\cdot141233$ gives the possible divisors
below $p$ as $1,3,5,9,15,45$, none in that class.  For completeness,
$141233-1=2^4\cdot7\cdot13\cdot97$ and the Lucas test with base $3$ has
\[
 3^{141232}\equiv1\pmod{141233},
\]
while the residues at exponents $141232/q$, for $q=2,7,13,97$, are
$141232,85064,78498,79135$; each residue minus one is coprime to $141233$.
Thus $141233$ is prime and the divisor list is complete.  This antecedent
extremal failure remains an explicit negative control \cite{Localization}.

\section{Status, provenance and replay boundary}
This continuation gives a complete exterior counterpart to the previous middle
atlas. It does not close the unbounded existence step. The three original
sources in Theorem~\ref{thm:atlas} may still, logically, all be empty at a
prescribed prime. A proof of a nonzero original coefficient, or a genuine
prime-parameter cutoff with its entire exceptional domain certified, is still
required before treating the proposed Turn~8 existence algorithm as established.

The main standard-library verifier compares every original first-half exterior
state with the three-chart atlas at every $p\equiv1\pmod4$ prime through its
recorded bound. The separate checker imports neither it nor predecessor code
and enumerates primitive $h,r,s$ directly. Both verify the hard example (14),
the source-domain counterexamples, and the original integer family (15).
The supplied execution receipt records actual outcomes, not outcomes inferred
from the existence of source files. Replaying previous Turn~6--7 files repairs
the downloadable handoff; it is not a new independent review of their proofs.
No universal ES theorem, whole-history audit, Lean build, new overall ES
verification range or historical-priority determination is claimed.
